% ===========================================================================
%  10 — Telemetria
% ===========================================================================
\chapter{Telemetria}
\label{cap:telemetria}

\begin{decisao}
\textbf{Dois canais de telemetria, não um.} A plataforma tem dois fluxos de
evento com garantias, latências e destinos diferentes, e confundi-los quebra um
dos dois:

\begin{itemize}
  \item o canal \textbf{operacional} (motor), \emph{síncrono e transacional},
    lido em tempo quase real para curar o próximo bloco de \emph{feed};
  \item o canal \textbf{analítico} (SDK), \emph{assíncrono e em lote}, para
    análise posterior.
\end{itemize}
\end{decisao}

\section{Por que dois canais}

\begin{tabularx}{\linewidth}{@{}L{26mm}L{34mm}L{34mm}@{}}
\toprule
\headrow \thd{Aspecto} & \thd{Operacional (motor)} & \thd{Analítico (SDK)}\\
\midrule
Rota & \rota{POST /runs/:id/eventos},
  \rota{GET /runs/:id/feed} &
  \rota{POST /telemetry/batch}, \rota{/beacon}, \rota{/ingest}\\
\zebra Transporte & HTTP direto, dentro da transação &
  HTTP \seta{} Pub/Sub \seta{} \emph{push} \seta{} HTTP\\
Latência & Milissegundos; a próxima tela depende dela &
  Segundos a minutos; ninguém espera\\
\zebra Ordenação & \textbf{Garantida} pela transação e pelo \emph{lock} da
  \cd{run} & Sem garantia\\
Destino & \cd{exp08ef1b\_events} & \cd{telemetria\_eventos} (particionada)\\
\zebra Entrega & Exatamente uma vez (transacional) &
  \emph{At least once}; \emph{dedupe} por \cd{(event\_id, ts)}\\
Quem lê & O motor de curadoria, o painel do professor, o dossiê &
  Ninguém, hoje --- é dado de análise futura\\
\bottomrule
\end{tabularx}

\begin{regra}
\textbf{\emph{Dwell} NUNCA passa pelo Pub/Sub.} A latência e a falta de ordenação
quebrariam a adaptação do \emph{feed} e o \emph{lock} por \cd{run}: o motor
precisa saber, no instante em que cura o próximo bloco, exatamente quais posts o
estudante viu e por quanto tempo. Um evento que chega dois minutos depois, fora de
ordem, produziria uma curadoria baseada em passado errado --- e o estudante não teria
como perceber.
\end{regra}

\begin{nota}
A recíproca também vale, e é o que justifica o segundo canal: mandar \emph{tudo}
pelo caminho transacional colocaria, no pico de quatro mil simultâneos, um
\cd{INSERT} sincronizado por interação de estudante dentro da mesma transação que
serve a tela. O canal analítico existe para absorver volume sem competir com o
que o estudante está esperando.
\end{nota}

\section{O canal analítico, ponta a ponta}

\begin{figure}[H]
\centering
\fitpage{%
\begin{tikzpicture}[x=1cm, y=1cm]

  % Três linhas: o caminho normal (em cima), o beacon do fechamento (meio) e a
  % ingestão (embaixo). As duas descidas até o Pub/Sub e a partir dele correm
  % paralelas, sem se cruzar.
  \node[etapaio, text width=22mm, minimum height=15mm] (exp) at (1.35,0)
    {\textbf{Experiência}\desc{\texttt{track(verbo,}\\\texttt{objeto, dados)}}};

  \node[etapa, text width=28mm, minimum height=15mm] (sdk) at (4.95,0)
    {\textbf{\texttt{TelemetryClient}}\desc{\emph{buffer} em memória;\\
     lote de 50 ou \SI{10}{\second}}};

  \node[etapa, text width=28mm, minimum height=15mm] (batch) at (9.45,0)
    {\texttt{POST}\\\texttt{/telemetry/batch}\desc{\texttt{OidcGuard};\\identidade vem do TOKEN}};

  \node[fita, text width=24mm, minimum height=14mm] (topic) at (14.3,0)
    {\textbf{Pub/Sub}\desc{\emph{topic} \texttt{-telemetria}}};

  \node[etapaerr, text width=28mm, minimum height=13mm] (beacon) at (9.45,-3.2)
    {\texttt{POST}\\\texttt{/telemetry/beacon}\desc{\texttt{pagehide}; \emph{token} NO\\CORPO; 204 sempre}};

  \node[etapa, text width=36mm, minimum height=15mm] (ingest) at (11.4,-6.3)
    {\texttt{POST}\\\texttt{/telemetry/ingest}\desc{\texttt{PubSubPushGuard}:\\ID \emph{token} do Google,\\\texttt{aud} e \texttt{email}}};

  \node[etapaout, text width=42mm, minimum height=16mm] (pg) at (5.6,-6.3)
    {\textbf{\texttt{telemetria\_eventos}}\desc{\texttt{INSERT} multi-linha,\\
     \texttt{ON CONFLICT (event\_id, ts)}\\\texttt{DO NOTHING};\\particionada por mês}};

  \draw[pipe] (exp) -- (sdk);
  \draw[pipe] (sdk) -- (batch) node[pipelbl, midway, above=1pt]{$\leq$ 50\\eventos};
  \draw[pipe] (batch) -- (topic);
  \draw[pipe] (sdk.south) |- (beacon.west)
    node[pipelbl, pos=0.75, above=1pt]{\texttt{sendBeacon}};
  \draw[pipe] (beacon.east) -- (13.6,-3.2) -- (13.6,0 |- topic.south)
    node[pipelbl, pos=0.5, left=2pt, align=right]{publica no\\mesmo \emph{topic}};
  \draw[pipe] (15.0,0 |- topic.south) |- (ingest.east)
    node[pipelbl, pos=0.3, left=2pt, align=right]{\emph{push}\\autenticado};
  \draw[pipe] (ingest) -- (pg);

  \node[actnote, text width=70mm, anchor=north west] at (0.05,-8.3)
    {\textbf{Onde há perda possível, e só ali:} entre o \texttt{track()} e o
     \emph{flush}. Se o processo do navegador morrer com o \emph{buffer} cheio ---
     aba fechada à força, queda de energia --- aqueles eventos somem. O
     \texttt{sendBeacon} do \texttt{pagehide} cobre o fechamento normal. A partir
     do Pub/Sub, a entrega é \textbf{durável}: é exatamente por isso que um
     \emph{buffer} em Redis não substituiria o Pub/Sub --- ele reintroduziria a
     perda que o Pub/Sub eliminou.};

  \node[actnote, text width=70mm, anchor=north west] at (7.95,-8.3)
    {\textbf{Por que o \emph{beacon} responde 204 sempre:} o navegador já
     descartou a página e ninguém lê o corpo. Erro aqui só serviria para poluir o
     log do cliente. Falha de verificação de \emph{token} vira descarte
     silencioso --- com aviso no servidor, não no cliente. O \emph{token} vai no
     \textbf{corpo} porque \texttt{sendBeacon} não envia cabeçalho, e isso é
     limitação da API do navegador, não escolha.};

\end{tikzpicture}}
\begin{legendadiag}
\amostra{etapaio} origem (o código da experiência) \quad
\amostra{etapa} caminho normal \quad
\amostra{etapaerr} caminho do fechamento da página \quad
\amostra{etapaout} destino durável
\end{legendadiag}
\caption[Canal analítico de telemetria]{O canal analítico ponta a ponta. Os
  únicos eventos que podem se perder são os que estão no \emph{buffer} em memória
  do navegador no instante em que o processo morre; do Pub/Sub em diante a
  entrega é durável.}
\label{fig:telemetria-analitica}
\end{figure}

\section{A rota de ingestão e sua tranca}

\rota{POST /telemetry/ingest} é a rota mais delicada da API: ela grava dado de
estudante, o \cd{escolaId} vem no corpo (já enriquecido pelo \rota{/batch}) e ela
\textbf{não passa pelo \cd{OidcGuard}} --- porque quem chama é o Pub/Sub, não um
usuário, e não existe \emph{token} do hub nessa chamada.

\begin{regra}
Sem \cd{PubSubPushGuard}, a rota seria um \emph{endpoint} público que escreve
telemetria falsa em qualquer escola. O \emph{guard} verifica \textbf{três} coisas
no ID token que o Pub/Sub assina:

\begin{enumerate}
  \item \textbf{emissor Google} (\cd{accounts.google.com}), pelo JWKS público;
  \item \cd{aud} igual a \env{TELEMETRIA\_PUSH\_AUDIENCE};
  \item \cd{email} igual a \env{TELEMETRIA\_PUSH\_SA}.
\end{enumerate}

As três juntas é que fazem a prova. A assinatura sozinha só diz \enquote{algum
projeto GCP}, que não é ninguém em particular.

\textbf{Falha fechada:} sem as duas variáveis, a rota rejeita tudo com 401 ---
nunca fica aberta. O Terraform casa os dois lados: a \emph{subscription} é criada
com \cd{oidc\_token} dessa mesma conta de serviço e audiência.
\end{regra}

\begin{decisao}
\textbf{\emph{Push subscription} para a própria API, e não um \emph{worker}
dedicado.} O desenho original previa
\cd{\{contrato\}-telemetry-worker} como serviço próprio. A \emph{push
subscription} entrega o mesmo resultado com uma peça a menos para operar --- um
serviço a menos com \emph{deploy}, log, escala e \emph{alerting} próprios. O
\emph{worker} separado volta se o volume justificar isolar o consumo do tráfego
de estudante.
\end{decisao}

\section{Destino: PostgreSQL, não BigQuery}

\begin{decisao}
\textbf{O canal analítico grava no mesmo PostgreSQL} (decisão de 2026-08-27,
\adrr{ADR-14}).

Avaliou-se mandar o canal analítico para o BigQuery por \emph{subscription}
nativa do Pub/Sub. Não foi, por dois motivos concretos: (a) os relatórios do
professor \textbf{não leem} \cd{telemetria\_eventos} --- eles leem
\cd{exp08ef1b\_events}, o canal síncrono do motor ---, então o BigQuery não
resolvia problema nenhum \emph{hoje}; e (b) criava uma segunda plataforma de
dados para operar, com IAM, custo e ferramentas próprios.

Ficando no mesmo banco, se um relatório precisar desse dado um dia, é um
\cd{JOIN} local com \cd{turmas} e \cd{users} --- e é uma plataforma de dados a
menos.
\end{decisao}

O que sustenta essa escolha na escala alvo são duas medidas, e nenhuma delas é um
serviço novo:

\subsection{\cd{INSERT} multi-linha}

\begin{lstlisting}[language=tsunoo, caption={O gargalo real, e a correção (\arq{src/telemetry/telemetry.controller.ts})}, label={lst:insert}]
// UM INSERT multi-linha, nao um por evento. E o gargalo real deste caminho:
// um lote de 50 eventos custava 50 round-trips ao banco, e no pico da manha
// (~4k estudantes simultaneos) isso e ~800 INSERTs/s numa instancia pequena.
// Em lote vira ~16/s, que e trivial.
//
// Nao e caso para Redis: quem amortece o pico ja e o Pub/Sub, que ainda por
// cima e DURAVEL — um buffer em memoria perderia eventos no processo que
// morre antes do flush.
const linhas = doGrupo.map(
  (e) => Prisma.sql`(
    ${e.eventId}::uuid, ${new Date(e.ts)}, ${e.escolaId}, ${e.estudanteId}, ${e.expId},
    ${e.expVersion}, ${e.sessaoId}::uuid, ${e.verbo}, ${e.objeto},
    ${JSON.stringify(e.resultado ?? null)}::jsonb,
    ${JSON.stringify(e.contexto  ?? null)}::jsonb)`,
);
await tx.$executeRaw(Prisma.sql`
  INSERT INTO telemetria_eventos
    (event_id, ts, escola_id, estudante_id, exp_id, exp_version, sessao_id,
     verbo, objeto, resultado, contexto)
  VALUES ${Prisma.join(linhas)}
  ON CONFLICT (event_id, ts) DO NOTHING`);
\end{lstlisting}

\begin{nota}
O lote é \textbf{agrupado por escola} antes da gravação, e cada grupo grava na
sua própria \emph{tenancy} (\cd{comTenancy(fn, escolaId)}). Um lote do Pub/Sub
pode conter eventos de escolas diferentes, e uma única transação não poderia ter
duas políticas de RLS armadas ao mesmo tempo.
\end{nota}

\subsection{Partição mensal}

Detalhada na Seção~\ref{sec:particao}. Em resumo: \cd{PARTITION BY RANGE (ts)},
partições no \emph{schema} \cd{particoes} (fora do alcance do Prisma), chave
primária composta \cd{(event\_id, ts)} porque o PostgreSQL exige a chave de
partição em todo índice único, partição \cd{DEFAULT} como rede de segurança, e
faixas declaradas até 2028-12.

\begin{decisao}
\textbf{Nada de Redis como \emph{buffer}} (\adrr{ADR-14}). Quem amortece o pico
já é o Pub/Sub --- e ele é \emph{durável}. Um \emph{buffer} em memória perde os
eventos do processo que morre antes do \emph{flush}, ou seja, reintroduz
exatamente a perda que o Pub/Sub eliminou. Redis só entra com gargalo medido, e o
gargalo medido aqui era o laço de \cd{INSERT}.
\end{decisao}

\section{Privacidade da telemetria}

\begin{regra}
\begin{itemize}
  \item O evento carrega \cd{estudante\_id} \textbf{opaco} (UUID), nunca nome ou
    \emph{e-mail}.
  \item A identidade é \textbf{derivada do \emph{token}} no servidor. O que o
    cliente enviou nesses campos é descartado --- então o cliente não
    \emph{pode} colocar PII ali.
  \item Nenhuma PII em log em texto claro.
  \item Texto livre sensível (o \emph{check-in} de bem-estar) \textbf{nunca} vai
    à telemetria; só contagens.
  \item Em \cd{app\_70\_sessoes.payload}, \cd{eventos[].dados} aceita apenas
    número, booleano ou \emph{string} de um \textbf{vocabulário fechado}. O campo
    de texto livre que existia era o único canal de texto arbitrário do
    \emph{payload} --- e o caminho por onde o nome de uma criança chegaria a um
    JSONB de longa duração. Trava estrutural, não lista de proibidos.
\end{itemize}
\end{regra}

\begin{limitacao}
\textbf{Retenção sem prazo declarado} também aqui. \cd{telemetria\_eventos}
guarda evento por \cd{estudante\_id} de \emph{todas} as experiências e não tem prazo
--- pela mesma razão registrada na Seção~\ref{sec:particao} e em \adrr{ADR-13}:
prazo só existe junto com a rotina que o cumpre. A diferença é que, particionada,
essa rotina passou a ser trivial de escrever (\cd{DROP TABLE
particoes.<mês>}), o que remove o último argumento técnico para não ter prazo.
A pergunta que sobra é de produto e de contrato, e é da plataforma inteira.
\end{limitacao}

\section{Comportamento em desenvolvimento}

\begin{tabularx}{\linewidth}{@{}L{40mm}X@{}}
\toprule
\headrow \thd{Configuração} & \thd{Comportamento}\\
\midrule
\env{TELEMETRIA\_TOPIC} vazio & Os eventos são apenas \emph{logados} no console
  (\cd{[telemetria/dev] N eventos}). Não há Pub/Sub em desenvolvimento e nada
  chega a \rota{/ingest}.\\
\zebra \env{TELEMETRIA\_PUSH\_SA} / \env{..\_AUDIENCE} vazias &
  \rota{/telemetry/ingest} responde 401 a tudo. É o comportamento correto: em
  desenvolvimento ninguém deveria chamá-la.\\
\bottomrule
\end{tabularx}

\begin{nota}
Essa combinação é o que permite rodar a plataforma inteira, com telemetria
funcionando de ponta a ponta do lado do cliente, sem nenhuma credencial de GCP
--- e é o que faz \rf{RNF-14} (testabilidade sem nuvem) ser verdade e não
aspiração.
\end{nota}
